%% kapa: exam=1 round=26 year=2015 period=2 form=A subject=accounting subject_ko=회계학 seq=15 number=55
%% source: Q-Net 감정평가사 기출문제 (cst00309) artlSeq=5138001 file=extracted/1008/5138001/제26회_감정평가사 1차_기출/제26회_감정평가사 1차_기출/2교시_감평_A형.pdf
%% conversion: kordoc 4.24.51 -> tools/gen1.py (auto); review=proofread-sonnet; flags=none
\begin{question}{55}
(주)감평은 20×3년 초 토지를 ￦1,500,000에 취득하고 매년 말 공정가치로 평가하는 재평가모형을 적용한다. 또한 재평가잉여금을 자산의 처분시점에 이익잉여금으로 직접 대체하기로 하였다. 동 토지의 매년 말 공정가치는 다음과 같다.
\begin{framedtext}
\centering\begin{tabular}{cc}
\uline{20×3년 말} & \uline{20×4년 말} \\
￦1,200,000 & ￦1,600,000
\end{tabular}\par
\end{framedtext}
(주)감평이 20×5년 말에 동 토지를 ￦1,100,000에 처분했을 때, 토지의 보유 및 처분과 관련하여 다음의 설명 중 옳지 않은 것은?
\begin{choices}
  \item 20×3년 초부터 20×5년 말까지 이익잉여금이 총 ￦400,000 감소한다.
  \item 20×3년 당기순이익이 ￦300,000 감소한다.
  \item 20×4년 당기순이익이 ￦200,000 증가한다.
  \item 20×4년 기타포괄이익이 ￦100,000 증가한다.
  \item 20×5년 유형자산처분손실이 ￦500,000 인식된다.
\end{choices}
\end{question}
\answer{3}
